% ================================================================================
% Paquetes y comandos propios del trabajo
% ================================================================================
\usepackage{tikz}
\usetikzlibrary{positioning, arrows.meta, shapes.geometric}
\usepackage{url}
\urlstyle{same}

% Recuadro para marcar lo que todavía falta medir o completar.
% Se busca en el PDF por la palabra "Pendiente".
\newcommand{\pendiente}[1]{%
	\par\noindent
	\fcolorbox{orange}{yellow!15}{%
		\parbox{\dimexpr\linewidth-2\fboxsep-2\fboxrule\relax}{%
			\small\textbf{Pendiente:} #1}}%
	\par}

% Espacio reservado para una figura que se agregará después.
\newcommand{\figurapendiente}[2]{%
	\begin{figure}[H]
		\centering
		\fbox{\parbox[c][5cm][c]{0.8\linewidth}{\centering\small
			Espacio reservado para la figura\\[0.5ex] \textit{#1}}}
		\caption{#1}
		\label{#2}
	\end{figure}}

% Notación de las métricas
\newcommand{\tscan}{\ensuremath{T_{\mathit{scan}}}}
\newcommand{\tfuno}{\ensuremath{T_{\mathit{F1}}}}
\newcommand{\ttotal}{\ensuremath{T_{\mathit{total}}}}
\newcommand{\tlearn}{\ensuremath{T_{\mathit{learn}}}}
\newcommand{\tstart}{\ensuremath{T_{\mathit{start}}}}
